\chapter{Problemas originales de consolidación preexamen}
\label{ch:problemas-preexamen}
\noindent Estos ocho ejercicios complementarios practican los métodos presentes en los apuntes aportados; las cifras y formulaciones se han creado para esta edición. No sustituyen los problemas del profesor ni son datos de equipos reales. Resuelva sin mirar las respuestas. En todos los casos, identifique fórmula, hipótesis, unidades y limitaciones.

\section*{PE01. Fractura frágil de placa con grieta central}
Una placa ideal muy ancha presenta una grieta pasante central cuya longitud total es $2a=6\,\mathrm{mm}$. Adopte $\sigma=200\,\mathrm{MPa}$, $Y=1.12$ y una tenacidad aplicable supuesta de $K_{\mathrm{mat}}=38\,\mathrm{MPa\sqrt m}$. Calcule $K_I$, la semilongitud crítica $a_c$ y la longitud total crítica. Indique qué hipótesis de validez tendría que examinar para usar el criterio sobre una pieza real.

\section*{PE02. Interpolación en curva S--N}
Para una idealización en régimen de Basquin se conocen dos puntos $(N_1,\sigma_{a1})=(10^3,500\,\mathrm{MPa})$ y $(N_2,\sigma_{a2})=(10^6,250\,\mathrm{MPa})$. Suponga una ley $\sigma_a=A N^b$ entre ambos valores. Obtenga $b$ y prediga $\sigma_a$ para $N=10^5$ ciclos. Comente el efecto de la tensión media y los límites de extrapolación.

\section*{PE03. Integración de la ley de Paris}
Para un ejemplo ficticio se adopta $da/dN=C(\Delta K)^m$, con $m=3$, $C=2\times10^{-12}\,\mathrm{m/ciclo}/(\mathrm{MPa\sqrt m})^3$, $Y=1.1$, $\Delta\sigma=100\,\mathrm{MPa}$, $a_i=1\,\mathrm{mm}$ y $a_f=5\,\mathrm{mm}$. Integre con $\Delta K=Y\Delta\sigma\sqrt{\pi a}$ y calcule el número idealizado de ciclos. Explique por qué no representa una vida garantizada.

\section*{PE04. Sensibilidad de la fluencia estacionaria}
La tasa de fluencia secundaria idealizada es $10^{-9}\,\mathrm{s^{-1}}$ a $100\,\mathrm{MPa}$. Con temperatura constante y exponente $n=4$, estime el factor de aumento de tasa a $120\,\mathrm{MPa}$ y el tiempo necesario para acumular una deformación de fluencia de $0.005$ en los dos casos. Explique por qué no equivale al tiempo de rotura.

\section*{PE05. Barra con dilatación térmica impedida}
Una barra de $E=110\,\mathrm{GPa}$, $\alpha=9\times10^{-6}\,\mathrm{K^{-1}}$ y calentamiento uniforme $\Delta T=80\,\mathrm K$ está completamente impedida de alargarse. Calcule la tensión elástica térmica con signo y determine si el modelo sería admisible si el límite elástico a esa temperatura fuese $60\,\mathrm{MPa}$.

\section*{PE06. Eco ultrasónico}
En un ensayo ideal por pulso-eco se usa velocidad longitudinal de $5900\,\mathrm{m/s}$ y un eco de reflexión a $5.0\,\mu\mathrm s$. Determine la profundidad equivalente del reflector, suponiendo trayectoria normal de ida y vuelta. Enumere tres causas por las que no puede identificarse directamente con el tamaño real del defecto.

\section*{PE07. Dos criterios de vencimiento}
En un programa completamente ficticio una tarea debe repetirse cada $100\,\mathrm{FH}$ o $6$ meses, lo que ocurra primero. Última ejecución: 1 de abril de 2026 a $3900\,\mathrm{FH}$. Estado: 1 de junio de 2026 a $3960\,\mathrm{FH}$. Calcule el próximo umbral de horas, fecha de calendario y márgenes. Si se vuelan $5\,\mathrm{FH}$ por jornada de uso, estime cuántas jornadas quedan por horas. Indique qué evidencia de cumplimiento faltaría.

\section*{PE08. Fiabilidad y tasa constante}
Una población académica acumula $20\,000\,\mathrm{FH}$ con 4 retiradas observadas. En un modelo de Poisson homogéneo y tasa constante, estime la tasa por hora, el tiempo medio entre retiradas y la probabilidad de ninguna retirada en las próximas $200\,\mathrm{FH}$. Explique por qué estos resultados no determinan la vida límite aprobada.
